\chapter{Исследовательская часть}

\section{Модель вычислений}
Для оценки трудоёмкости используется следующая модель вычислений.
\begin{enumerate}
	\item Операции \texttt{=}, \texttt{+}, \texttt{-} (в том числе унарный минус), \texttt{+=}, \texttt{-=}, \texttt{==}, \texttt{!=}, \texttt{<}, \texttt{>}, \texttt{<=}, \texttt{>=}, \texttt{[]}, \texttt{++}, \texttt{--}, \texttt{<<}, \texttt{>>}, \texttt{\&\&}, \texttt{||} имеют трудоёмкость 1.
	\item Операции \texttt{*}, \texttt{/}, \texttt{\%}, \texttt{*=}, \texttt{/=}, \texttt{\%=} имеют трудоёмкость 2.
	\item Трудоёмкость условного оператора
	\begin{equation}
		f_{if} = f_{\text{усл}} +
		\begin{cases}
			\min(f_A, f_B) & \text{в лучшем случае}, \\
			\max(f_A, f_B) & \text{в худшем случае},
		\end{cases}
	\end{equation}
	где $f_{\text{усл}}$~--- трудоёмкость вычисления условия, $f_A$ и $f_B$~--- трудоёмкости ветвей.
	\item Трудоёмкость цикла
	\begin{equation}
		\label{eq:f-for}
		f_{for} = f_{\text{иниц}} + f_{\text{сравн}} + K \cdot (f_{\text{тела}} + f_{\text{инкр}} + f_{\text{сравн}}),
	\end{equation}
	где $K$~--- число итераций. Для цикла \texttt{for i in 0..<x} языка Swift диапазон \texttt{0..<x} создаётся один раз перед началом цикла, поэтому вычисление границы $x$ относится к инициализации: $f_{\text{иниц}} = 1 + f_{\text{гр}}$, где $f_{\text{гр}}$~--- трудоёмкость вычисления границы. Сравнение и инкремент имеют трудоёмкость 1. Тогда
	\begin{equation}
		\label{eq:f-for-swift}
		f_{for} = 2 + f_{\text{гр}} + K \cdot (2 + f_{\text{тела}}).
	\end{equation}
	\item Трудоёмкость вызова функции равна 0.
\end{enumerate}

Выделение памяти под результирующую матрицу и вспомогательные массивы, а также определение размеров матриц в оценке не учитываются, так как выполняются одинаково во всех алгоритмах. Далее обозначено $H = \floor{m/2}$.

\section{Трудоёмкость алгоритмов}

\subsection{Стандартный алгоритм умножения матриц}
Тело внутреннего цикла \texttt{c[i][j] = c[i][j] + a[i][k] * b[k][j]} содержит 8 операций индексации, одно умножение, одно сложение и одно присваивание, его трудоёмкость равна $8 + 2 + 1 + 1 = 12$. Границы циклов $n$, $q$, $m$ не требуют вычислений. По формуле~\eqref{eq:f-for-swift}
\begin{equation}
	\label{eq:f-std}
	f_{\text{станд}} = 2 + n \bigl(4 + q (4 + 14m)\bigr) = 14nmq + 4nq + 4n + 2.
\end{equation}

\subsection{Алгоритм Винограда}
Трудоёмкость алгоритма складывается из трудоёмкостей четырёх этапов:
\begin{equation}
	f_{\text{Вин}} = f_{mulH} + f_{mulV} + f_{\text{осн}} + f_{\text{неч}}.
\end{equation}

Граница циклов по индексу пары $m / 2$ имеет трудоёмкость $f_{\text{гр}} = 2$.

Тело цикла вычисления $mulH$ \texttt{mulH[i] = mulH[i] + a[i][2*j] * a[i][2*j + 1]} содержит 6 операций индексации, 3 умножения, 2 сложения и присваивание: $6 + 6 + 2 + 1 = 15$. Тогда
\begin{equation}
	f_{mulH} = 2 + n (2 + 4 + 17H) = 17nH + 6n + 2.
\end{equation}
Цикл вычисления $mulV$ устроен аналогично:
\begin{equation}
	f_{mulV} = 17qH + 6q + 2.
\end{equation}

В основном цикле присваивание \texttt{c[i][j] = -mulH[i] - mulV[j]} имеет трудоёмкость $2 + 1 + 1 + 1 + 1 + 1 = 7$ (4 индексации, унарный минус, вычитание, присваивание). Тело внутреннего цикла
\texttt{c[i][j] = c[i][j] + (a[i][2*k] + b[2*k + 1][j]) * (a[i][2*k + 1] + b[2*k][j])}
содержит 12 операций индексации, 5 умножений, 5 сложений и присваивание: $12 + 10 + 5 + 1 = 28$. Тогда
\begin{equation}
	f_{\text{осн}} = 2 + n \Bigl(4 + q \bigl(7 + 4 + 2 + 30H\bigr)\Bigr) = 30nqH + 13nq + 4n + 2.
\end{equation}

Проверка условия \texttt{m \% 2 == 1} имеет трудоёмкость 3, тело цикла \texttt{c[i][j] = c[i][j] + a[i][m - 1] * b[m - 1][j]}~--- $8 + 2 + 2 + 1 + 1 = 14$:
\begin{equation}
	f_{\text{неч}} = 3 +
	\begin{cases}
		0, & m \text{ чётное (лучший случай)}, \\
		16nq + 4n + 2, & m \text{ нечётное (худший случай)}.
	\end{cases}
\end{equation}

Итого
\begin{equation}
	\label{eq:f-wino}
	\begin{split}
	f_{\text{Вин}} = {} & 30nqH + 17H(n + q) + 13nq + 10n + 6q + 9 \; + \\
	& + \begin{cases}
		0, & m \text{ чётное}, \\
		16nq + 4n + 2, & m \text{ нечётное}.
	\end{cases}
	\end{split}
\end{equation}

\subsection{Оптимизированный алгоритм Винограда}
Оптимизации А и Б не меняют структуру алгоритма Винограда, поэтому трудоёмкость каждой реализации выражается через трудоёмкость базового алгоритма~\eqref{eq:f-wino} с учётом изменённых операций.

\subsubsection{Реализация с оптимизацией А}
Замена $m / 2$ на $m \gg 1$ уменьшает трудоёмкость вычисления границы цикла с 2 до 1. Так как граница вычисляется один раз при входе в цикл, выигрыш равен числу входов в циклы по индексу пары: $n$ раз при вычислении $mulH$, $q$ раз при вычислении $mulV$ и $nq$ раз в основном цикле:
\begin{equation}
	\label{eq:f-shift}
	f_{\text{А}} = f_{\text{Вин}} - (nq + n + q).
\end{equation}

Если бы граница проверялась в условии цикла на каждой итерации, как в цикле \texttt{for} языка~C, выигрыш составил бы $1$ на каждую итерацию, то есть $nqH + nH + qH$. В языке Swift такой эффект отсутствует.

\subsubsection{Реализация с оптимизацией Б}
Составной оператор \texttt{+=} избавляет от повторной индексации накапливаемого элемента и от отдельной операции сложения:
\begin{itemize}
	\item \texttt{mulH[i] += a[i][2*j] * a[i][2*j + 1]}: $1 + 1 + 4 + 5 + 2 = 13$ вместо 15;
	\item \texttt{c[i][j] += (a[i][2*k] + b[2*k + 1][j]) * (a[i][2*k + 1] + b[2*k][j])}: 25 вместо 28;
	\item \texttt{c[i][j] += a[i][m - 1] * b[m - 1][j]}: 11 вместо 14.
\end{itemize}
Тогда
\begin{equation}
	\label{eq:f-add}
	\begin{split}
	f_{\text{Б}} = {} & 27nqH + 15H(n + q) + 13nq + 10n + 6q + 9 \; + \\
	& + \begin{cases}
		0, & m \text{ чётное}, \\
		13nq + 4n + 2, & m \text{ нечётное}.
	\end{cases}
	\end{split}
\end{equation}

\subsubsection{Реализация с оптимизациями А и Б}
При совместном применении оптимизаций их выигрыши складываются, так как они затрагивают разные операции: оптимизация А~--- вычисление границы циклов по индексу пары, оптимизация Б~--- накопление сумм. Поэтому
\begin{equation}
	\label{eq:f-shift-add}
	f_{\text{АБ}} = f_{\text{Б}} - (nq + n + q).
\end{equation}

\subsection{Сводная оценка для квадратных матриц}
Для квадратных матриц $n = m = q = N$ подстановка $H = N/2$ при чётном $N$ и $H = (N - 1)/2$ при нечётном $N$ в формулы~\eqref{eq:f-std}--\eqref{eq:f-shift-add} даёт результаты, приведённые в таблице~\ref{tbl:complexity}. Полученные формулы были дополнительно проверены непосредственным подсчётом трудоёмкости по модели для $N = \overline{1, 29}$.

\begin{table}[H]
\caption{Трудоёмкость алгоритмов для квадратных матриц размера $N \times N$}
\label{tbl:complexity}
\centering
\small
\begin{tabular}{|l|c|c|}
\hline
\textbf{Алгоритм} & \textbf{Чётное $N$} & \textbf{Нечётное $N$} \\
\hline
Стандартный & \multicolumn{2}{c|}{$14N^3 + 4N^2 + 4N + 2$} \\
\hline
Виноград & $15N^3 + 30N^2 + 16N + 9$ & $15N^3 + 31N^2 + 3N + 11$ \\
\hline
Виноград, опт. А & $15N^3 + 29N^2 + 14N + 9$ & $15N^3 + 30N^2 + N + 11$ \\
\hline
Виноград, опт. Б & $13{,}5N^3 + 28N^2 + 16N + 9$ & $13{,}5N^3 + 27{,}5N^2 + 5N + 11$ \\
\hline
Виноград, опт. А и Б & $13{,}5N^3 + 27N^2 + 14N + 9$ & $13{,}5N^3 + 26{,}5N^2 + 3N + 11$ \\
\hline
\end{tabular}
\end{table}

Согласно модели, алгоритм Винограда без оптимизаций более трудоёмкий относительно стандартного алгоритма: вдвое меньшее число умножений элементов матриц не компенсирует вычисление индексов $2k$ и $2k+1$ (4 умножения на итерацию) и большее число операций индексации. Оптимизация А влияет только на слагаемые порядка $N^2$ и практически не меняет трудоёмкость. Оптимизация Б уменьшает коэффициент при старшем члене с 15 до 13,5, и только версии алгоритма Винограда с этой оптимизацией по модели оказываются менее трудоёмкими, чем стандартный алгоритм. Совместное применение оптимизаций А и Б дополнительно уменьшает трудоёмкость версии с оптимизацией Б лишь на слагаемые порядка $N^2$. Все пять алгоритмов имеют кубическую трудоёмкость.

\section{Характеристики устройства}
Исследования проводились на машине со следующими характеристиками:
\begin{itemize}
	\item ноутбук MacBook Air;
	\item процессор Apple M5, 10 ядер (4 производительных и 6 энергоэффективных);
	\item оперативная память: 16 ГБ;
	\item операционная система: macOS 27.0.1.
\end{itemize}

\section{Инструменты замера процессорного времени}
\label{sec:time-tools}
Для замера процессорного времени использовалась функция \texttt{clock\_gettime\_nsec\_np} системной библиотеки macOS с идентификатором часов \texttt{CLOCK\_THREAD\_CPUTIME\_ID}~\cite{clock-gettime}. Функция возвращает значение часов в наносекундах в виде 64-битного беззнакового целого, а часы \texttt{CLOCK\_THREAD\_CPUTIME\_ID} учитывают процессорное время (в режиме пользователя и в режиме ядра), затраченное вызывающим потоком. Таким образом, в замер не попадает время, в течение которого поток был вытеснен другими процессами.

Время выполнения одного алгоритма для одного размера определяется как среднее по 10 повторам умножения одних и тех же случайно сгенерированных матриц с элементами из отрезка $[-100; 100]$. Чтобы компилятор не исключил вычисление неиспользуемого результата, результат передаётся в функцию \texttt{blackHole}, запрещённую к встраиванию. Замеры выполнялись без оптимизаций компилятора: программа собиралась командой \texttt{swift run -c release -Xswiftc -Onone}, где \texttt{-Onone}~--- уровень оптимизации компилятора Swift, аналогичный флагу \texttt{-O0} компиляторов C и C++. Конфигурация \texttt{release} использована для того, чтобы исключить отладочный вывод алгоритмов (блоки \texttt{\#if DEBUG}), а флаг \texttt{-Onone}~--- чтобы компилятор не устранял различия между реализациями, которые оцениваются моделью вычислений: повторную индексацию, вычисление индексов $2k$ и $2k + 1$, запись накопления суммы. Функция замера \texttt{measure} и функция \texttt{runTimeMeasurements}, выполняющая массированные замеры, приведены в листинге~\ref{lst:measure} приложения~А.

\section{Результаты замеров}
Замеры выполнялись для квадратных матриц чётных размеров $N = 10, 20, \ldots, 200$ и нечётных размеров $N = 11, 21, \ldots, 201$. Результаты приведены в таблицах~\ref{tbl:time}, \ref{tbl:time-odd} и на рисунках~\ref{img:time-even} и~\ref{img:time-odd}.

\begin{table}[H]
\caption{Процессорное время умножения квадратных матриц чётного размера, мс}
\label{tbl:time}
\centering
\begin{tabular}{|c|c|c|c|c|c|}
\hline
$N$ & Станд. & Виногр. & Опт. А & Опт. Б & Опт. А и Б \\
\hline
10 & 0{,}33 & 0{,}22 & 0{,}19 & 0{,}21 & 0{,}15 \\ \hline
20 & 0{,}95 & 0{,}60 & 0{,}62 & 0{,}64 & 0{,}55 \\ \hline
30 & 3{,}03 & 2{,}01 & 1{,}98 & 1{,}77 & 1{,}80 \\ \hline
40 & 6{,}84 & 4{,}46 & 4{,}39 & 3{,}92 & 4{,}02 \\ \hline
50 & 13{,}07 & 8{,}31 & 8{,}29 & 7{,}43 & 7{,}52 \\ \hline
60 & 22{,}14 & 14{,}25 & 14{,}34 & 12{,}86 & 12{,}84 \\ \hline
70 & 34{,}82 & 22{,}32 & 22{,}55 & 19{,}82 & 20{,}00 \\ \hline
80 & 51{,}50 & 33{,}01 & 33{,}00 & 29{,}70 & 31{,}73 \\ \hline
90 & 75{,}37 & 47{,}89 & 46{,}88 & 45{,}23 & 42{,}13 \\ \hline
100 & 102{,}85 & 64{,}97 & 63{,}88 & 58{,}31 & 58{,}02 \\ \hline
110 & 134{,}42 & 86{,}93 & 86{,}47 & 77{,}11 & 80{,}30 \\ \hline
120 & 175{,}97 & 111{,}37 & 110{,}01 & 100{,}08 & 99{,}94 \\ \hline
130 & 229{,}09 & 142{,}44 & 140{,}62 & 129{,}02 & 129{,}95 \\ \hline
140 & 278{,}98 & 175{,}54 & 173{,}75 & 157{,}37 & 157{,}88 \\ \hline
150 & 350{,}77 & 218{,}67 & 217{,}31 & 195{,}05 & 197{,}40 \\ \hline
160 & 416{,}21 & 261{,}64 & 258{,}35 & 229{,}73 & 230{,}48 \\ \hline
170 & 501{,}82 & 325{,}63 & 317{,}99 & 282{,}03 & 279{,}23 \\ \hline
180 & 596{,}14 & 372{,}26 & 372{,}56 & 336{,}77 & 343{,}29 \\ \hline
190 & 697{,}54 & 440{,}23 & 437{,}36 & 393{,}86 & 395{,}61 \\ \hline
200 & 807{,}31 & 512{,}29 & 503{,}37 & 454{,}18 & 462{,}96 \\ \hline
\end{tabular}
\end{table}

\begin{table}[H]
\caption{Процессорное время умножения квадратных матриц нечётного размера, мс}
\label{tbl:time-odd}
\centering
\begin{tabular}{|c|c|c|c|c|c|}
\hline
$N$ & Станд. & Виногр. & Опт. А & Опт. Б & Опт. А и Б \\
\hline
11 & 0{,}15 & 0{,}12 & 0{,}12 & 0{,}11 & 0{,}11 \\ \hline
21 & 1{,}07 & 0{,}72 & 0{,}78 & 0{,}64 & 0{,}73 \\ \hline
31 & 3{,}28 & 2{,}16 & 2{,}17 & 1{,}98 & 1{,}95 \\ \hline
41 & 7{,}17 & 4{,}75 & 4{,}68 & 4{,}30 & 4{,}34 \\ \hline
51 & 13{,}70 & 8{,}96 & 8{,}90 & 8{,}02 & 8{,}15 \\ \hline
61 & 22{,}79 & 14{,}97 & 14{,}88 & 13{,}41 & 13{,}39 \\ \hline
71 & 36{,}21 & 23{,}31 & 23{,}00 & 20{,}77 & 20{,}87 \\ \hline
81 & 58{,}57 & 37{,}46 & 36{,}04 & 31{,}87 & 31{,}78 \\ \hline
91 & 76{,}51 & 49{,}52 & 49{,}55 & 44{,}25 & 44{,}29 \\ \hline
101 & 104{,}10 & 66{,}70 & 65{,}60 & 59{,}43 & 59{,}51 \\ \hline
111 & 145{,}77 & 93{,}87 & 94{,}45 & 81{,}73 & 80{,}63 \\ \hline
121 & 181{,}42 & 116{,}56 & 114{,}43 & 103{,}34 & 103{,}54 \\ \hline
131 & 231{,}63 & 145{,}14 & 143{,}13 & 129{,}36 & 134{,}48 \\ \hline
141 & 289{,}01 & 183{,}66 & 181{,}30 & 159{,}26 & 162{,}10 \\ \hline
151 & 356{,}65 & 220{,}50 & 222{,}54 & 197{,}89 & 197{,}72 \\ \hline
161 & 419{,}44 & 269{,}68 & 265{,}67 & 238{,}55 & 243{,}00 \\ \hline
171 & 502{,}62 & 319{,}60 & 322{,}41 & 286{,}97 & 290{,}85 \\ \hline
181 & 607{,}20 & 380{,}53 & 377{,}57 & 339{,}79 & 343{,}41 \\ \hline
191 & 713{,}42 & 454{,}41 & 446{,}40 & 404{,}07 & 403{,}54 \\ \hline
201 & 828{,}41 & 533{,}75 & 517{,}29 & 464{,}74 & 463{,}35 \\ \hline
\end{tabular}
\end{table}


% Графики строятся по CSV-файлам из lab-02/TimeTestResult
\newcommand{\timeplot}[1]{%
	\begin{tikzpicture}
	\begin{axis}[
		width=0.97\linewidth, height=0.62\linewidth,
		xlabel={Линейный размер матриц $N$},
		ylabel={Процессорное время, мс},
		grid=major, grid style={dashed, gray!40},
		legend pos=north west, legend cell align=left,
		legend style={font=\small},
		xmin=0, xmax=210, ymin=0,
		x filter/.expression={mod(\thisrow{size}, 2) == #1 ? x : nan},
		unbounded coords=discard,
		scaled y ticks=false,
		yticklabel style={/pgf/number format/.cd, fixed, use comma},
	]
	\addplot[thick, blue, mark=*, mark size=1.3pt] table[x=size, y expr=\thisrow{time_ns}/1e6, col sep=comma] {../TimeTestResult/classic.csv};
	\addlegendentry{Стандартный}
	\addplot[thick, red, mark=square*, mark size=1.3pt] table[x=size, y expr=\thisrow{time_ns}/1e6, col sep=comma] {../TimeTestResult/winograd.csv};
	\addlegendentry{Виноград}
	\addplot[thick, green!50!black, dashed, mark=triangle*, mark size=1.6pt] table[x=size, y expr=\thisrow{time_ns}/1e6, col sep=comma] {../TimeTestResult/winograd_shift.csv};
	\addlegendentry{Виноград, опт. А}
	\addplot[thick, orange, dotted, mark=diamond*, mark size=1.6pt] table[x=size, y expr=\thisrow{time_ns}/1e6, col sep=comma] {../TimeTestResult/winograd_add_assigned.csv};
	\addlegendentry{Виноград, опт. Б}
	\addplot[thick, violet, dashdotted, mark=pentagon*, mark size=1.6pt] table[x=size, y expr=\thisrow{time_ns}/1e6, col sep=comma] {../TimeTestResult/winograd_shift_add_assigned.csv};
	\addlegendentry{Виноград, опт. А и Б}
	\end{axis}
	\end{tikzpicture}%
}

\begin{figure}[H]
	\centering
	\timeplot{0}
	\caption{Зависимость процессорного времени от линейного размера матриц, чётные размеры}
	\label{img:time-even}
\end{figure}

\begin{figure}[H]
	\centering
	\timeplot{1}
	\caption{Зависимость процессорного времени от линейного размера матриц, нечётные размеры}
	\label{img:time-odd}
\end{figure}

\section{Сравнительный анализ}
Для сравнения алгоритмов используется относительное значение трудоёмкости и времени работы алгоритма $X$ по отношению к алгоритму $Y$:
\begin{equation}
	\label{eq:rel}
	f_{\text{отн}} = \frac{f_X(N)}{f_Y(N)}, \qquad t_{\text{отн}} = \frac{t_X(N)}{t_Y(N)},
\end{equation}
где $f_X(N)$, $f_Y(N)$~--- трудоёмкости алгоритмов по таблице~\ref{tbl:complexity}, $t_X(N)$, $t_Y(N)$~--- измеренное процессорное время по таблицам~\ref{tbl:time} и~\ref{tbl:time-odd}. Значение меньше единицы означает, что алгоритм $X$ эффективнее алгоритма $Y$. В таблице~\ref{tbl:comparison} приведены значения при $N = 200$ по отношению к стандартному алгоритму.

\begin{table}[H]
\caption{Сравнение алгоритмов со стандартным алгоритмом при $N = 200$}
\label{tbl:comparison}
\centering
\begin{tabular}{|l|c|c|c|c|}
\hline
\textbf{Алгоритм} & \textbf{Трудоёмкость} & $\boldsymbol{f_{\text{отн}}}$ & \textbf{Время, мс} & $\boldsymbol{t_{\text{отн}}}$ \\
\hline
Стандартный & $112\,160\,802$ & $1{,}000$ & $807{,}31$ & $1{,}000$ \\
\hline
Виноград & $121\,203\,209$ & $1{,}081$ & $512{,}29$ & $0{,}635$ \\
\hline
Виноград, опт. А & $121\,162\,809$ & $1{,}080$ & $503{,}37$ & $0{,}624$ \\
\hline
Виноград, опт. Б & $109\,123\,209$ & $0{,}973$ & $454{,}18$ & $0{,}563$ \\
\hline
Виноград, опт. А и Б & $109\,082\,809$ & $0{,}973$ & $462{,}96$ & $0{,}573$ \\
\hline
\end{tabular}
\end{table}

При увеличении $N$ вдвое, со 100 до 200, время работы каждого алгоритма возрастает в 7,8--8,0 раза, то есть примерно в $2^3 = 8$ раз, что соответствует кубической трудоёмкости.

По данным замеров все версии алгоритма Винограда работают быстрее стандартного алгоритма при всех исследованных размерах: при $N = 200$ базовая версия быстрее в $1 / 0{,}635 \approx 1{,}6$ раза, версия с оптимизацией Б~--- в $1 / 0{,}563 \approx 1{,}8$ раза, версия с оптимизациями А и Б~--- в $1 / 0{,}573 \approx 1{,}7$ раза. Модель вычислений, напротив, предсказывает, что алгоритм Винограда на 8\,\% трудоёмче стандартного ($f_{\text{отн}} = 1{,}081$).

Расхождение связано с тем, что в модели обращение к элементу массива \texttt{[]} имеет такую же трудоёмкость, как сложение. В реализации матрица хранится как массив массивов \texttt{[[Int]]}, поэтому обращение к элементу $c_{i,j}$ требует обращения сначала к строке $c_i$, а затем к элементу строки. При изменении элемента массива Swift дополнительно проверяет, не используется ли массив другими переменными, и при необходимости копирует его (механизм копирования при записи, copy-on-write). Поэтому чтение и особенно изменение элемента матрицы обходятся дороже сложения или умножения целых чисел. Число обращений к элементам матриц во внутреннем цикле приведено в таблице~\ref{tbl:access}.

\begin{table}[H]
\caption{Число обращений к элементам матриц во внутреннем цикле}
\label{tbl:access}
\centering
\begin{tabular}{|l|c|c|c|}
\hline
\textbf{Алгоритм} & \textbf{Чтений $a$, $b$} & \textbf{Чтений $c$} & \textbf{Изменений $c$} \\
\hline
Стандартный & $2N^3$ & $N^3$ & $N^3$ \\
\hline
Виноград & $2N^3$ & $N^3/2$ & $N^3/2$ \\
\hline
Виноград, опт. А & $2N^3$ & $N^3/2$ & $N^3/2$ \\
\hline
Виноград, опт. Б & $2N^3$ & --- & $N^3/2$ \\
\hline
Виноград, опт. А и Б & $2N^3$ & --- & $N^3/2$ \\
\hline
\end{tabular}
\end{table}

Алгоритм Винограда обрабатывает за одну итерацию внутреннего цикла два слагаемых суммы, поэтому читает и изменяет элемент $c_{i,j}$ вдвое реже, чем стандартный алгоритм. Именно за счёт этого он работает быстрее, несмотря на дополнительные операции вычисления индексов $2k$ и $2k + 1$.

Оптимизации сравниваются с базовой версией алгоритма Винограда по формуле~\eqref{eq:rel}, где $Y$~--- алгоритм Винограда без оптимизаций. При $N = 200$:
\begin{itemize}
	\item для оптимизации А $f_{\text{отн}} = 121\,162\,809 / 121\,203\,209 \approx 1{,}000$ и $t_{\text{отн}} = 503{,}37 / 512{,}29 \approx 0{,}983$~--- оптимизация практически не влияет на трудоёмкость и время работы, так как граница цикла \texttt{0..<(m / 2)} вычисляется один раз при входе в цикл, а не на каждой итерации;
	\item для оптимизации Б $f_{\text{отн}} = 109\,123\,209 / 121\,203\,209 \approx 0{,}900$ и $t_{\text{отн}} = 454{,}18 / 512{,}29 \approx 0{,}887$~--- оптимизация уменьшает время работы на 11\,\%, что близко к предсказанным моделью 10\,\%: оператор \texttt{+=} изменяет элемент $c_{i,j}$ без отдельного чтения в правой части присваивания;
	\item для оптимизаций А и Б $f_{\text{отн}} = 109\,082\,809 / 121\,203\,209 \approx 0{,}900$ и $t_{\text{отн}} = 462{,}96 / 512{,}29 \approx 0{,}904$~--- выигрыш такой же, как от одной оптимизации Б. По сравнению с версией с оптимизацией Б ($Y$~--- версия с оптимизацией Б) $f_{\text{отн}} = 109\,082\,809 / 109\,123\,209 \approx 1{,}000$ и $t_{\text{отн}} = 462{,}96 / 454{,}18 \approx 1{,}019$: добавление оптимизации А выигрыша не даёт, а различие в 2\,\% находится в пределах погрешности замеров.
\end{itemize}

Для нечётного размера $N = 201$ значения $t_{\text{отн}}$ по отношению к стандартному алгоритму составляют $0{,}644$ для алгоритма Винограда, $0{,}624$ для версии с оптимизацией А, $0{,}561$ для версии с оптимизацией Б и $0{,}559$ для версии с оптимизациями А и Б, то есть близки к значениям для чётного размера. Дополнительный этап алгоритма Винограда для нечётного $m$ имеет трудоёмкость порядка $N^2$ и на соотношение алгоритмов не влияет.

\section{Вывод}
По результатам сравнительного анализа можно сделать следующие выводы:
\begin{itemize}
	\item все рассмотренные алгоритмы имеют кубическую трудоёмкость: при увеличении $N$ вдвое время работы возрастает примерно в 8 раз;
	\item все версии алгоритма Винограда работают быстрее стандартного алгоритма: при $N = 200$ базовая версия~--- в 1,6 раза, версии с оптимизацией Б и с оптимизациями А и Б~--- в 1,7--1,8 раза, так как алгоритм Винограда вдвое реже обращается к накапливаемому элементу результирующей матрицы;
	\item модель вычислений, в которой обращение к элементу массива имеет такую же трудоёмкость, как сложение, не предсказывает преимущество алгоритма Винограда над стандартным алгоритмом ($15N^3$ против $14N^3$ в старшем члене);
	\item оптимизация А не даёт выигрыша, что согласуется с моделью: граница цикла вычисляется один раз при входе в цикл;
	\item оптимизация Б уменьшает время работы алгоритма Винограда на 11\,\%, что близко к предсказанным моделью 10\,\%;
	\item совместное применение оптимизаций А и Б даёт такой же выигрыш, как одна оптимизация Б, так как вклад оптимизации А пренебрежимо мал;
	\item различий в соотношении алгоритмов для чётных и нечётных размеров матриц не обнаружено;
	\item из рассмотренных реализаций для умножения матриц предпочтителен алгоритм Винограда с оптимизацией Б (с оптимизацией А или без неё); алгоритм Винограда требует дополнительной памяти под массивы $mulH$ и $mulV$ размера $n + q$.
\end{itemize}
